\section{Whether the pilot argument warrants a larger experiment}
\label{exec13:dependent-design-review}

A valid conservative construction can be a poor next investment if it
cannot deliver a useful answer at the intended information level. The
independent-pilot argument therefore needs a precision review before a
large simulation campaign. The review preserves the earlier Gaussian
panel failures and evaluates the narrower bounded model on its own
terms. It does not transform known finite dependence into a guarantee
for unbounded Gaussian AR observations~\citep{exec13_ivrd_dependent}.

Use the authored sign-law example with independent fair signs
$X,V,H,E$,
\[
 Z=.3X+V,\qquad A=.8Z+.2X+H,\qquad Y=1.2A+.4X+E.
\]
It has $B_Z=13/10$, $B_A=56/25$, $B_Y=511/125$, $Q=4/5$ and
$M=24/25$. Grant an oracle its exact nuisance means, so the bias
radii are zero, and hold the empirical moments at their population
centers. With evaluation error allocation $.03$, the displayed
bounded-mean construction at $n=640$ contains zero in both moment
coordinates for each planning dependence range $d=0,4,19$.
The projected effect set is the whole line. These are deterministic
planning calculations at a fixed center, not prospective probabilities
or fitted-study outcomes.

For independent evaluation the same centered rectangle first has
a positive denominator lower endpoint at 2,075 rows. A more demanding
decision example stipulates an action threshold of 1.15, only .05
below the true coefficient. A sufficient centered condition for
the whole effect set to exceed that threshold is
\[
 8B_Z(B_Y+1.15B_A)
 \sqrt{\frac{\log(4(d+1)/.03)}{2\lfloor n/(d+1)\rfloor}}
 <M-1.15Q=.04.
\]
For $d=0$, the first integer sample size satisfying this displayed
bound is 7,344,273. This is not a proposed bank sample size. It
shows how the global range bound can obscure a small decision
margin even before fitting error is introduced. A sharper justified
variance-sensitive construction deserves development before a large
final campaign~\citep{exec13_ivrd_dependent}.

The alternative loss calibration can be conservative for another
reason. Even with exact conditional means and infinitely many
calibration observations, a zero lower bound on irreducible noise
leaves squared-loss limits $1$, $41/25$ and $2101/625$ in this law.
The resulting denominator bias radius is $\sqrt{41/25}>.8$ and
the numerator radius also exceeds $.96$. The centered rectangle
therefore remains uninformative as evaluation noise disappears.
These are certificate floors, not actual oracle fitting errors.
Known residual noise in a generated design can supply the missing
subtraction; estimated bank training loss cannot certify it by itself.

Qualification-screen power is separate from these precision calculations.
For the preserved 512 requests and sixteen-comparison family, at
least 500 successes are needed for a one-sided exact lower coverage
bound of .95. Exact binomial calculations give per-cell pass
probabilities about $.001849$, $.234929$, $.770054$ and $.997670$
when actual coverage is $.95$, $.97$, $.98$ and $.99$ respectively.
Thus a coverage guarantee of at least .95 alone does not promise
a high probability of passing the screen. The .99 finite-availability
rule is a raw fraction, so 507/512 suffices; a lower confidence
bound of .99 would be a different rule. The old high-dimension,
larger-support and nonlinear failures retain their original status
\citep{exec13_ivrd_dependent}.

The justified next work is a separately labeled bounded prototype:
fix a finite-dependent generator and independent pilot/calibration
law; implement and certify the actual learner; carry approximation,
numerical and selection errors into the bound; derive a sharper
evaluation region; and compare complete workflows at matched
information and resource levels. Strong, weak and zero relevance,
missing certificates, nonlinear omission and falsely declared
dependence all belong in the design. Only after development profiling
should fresh draws, replication budget, stopping rule, useful-width
scenario and certification-screen power be frozen. Unknown
dependence, arbitrary learners and causal instrument truth still
require new arguments or additional information. They do not
prevent the bounded proposal work, and no final qualification
study is approved by this feasibility calculation.
\label{exec13:dependent-design-review-end}
